\chapter{Lezione 9 - 22/10}

\section{Applicazioni Lineari}

\Definizione{Applicazione Lineare (Omomorfismo)}{
    Siano \(V\) e \(W\) due spazi vettoriali definiti sullo stesso campo \(K\).
    
    Un'applicazione \(T: V \to W\) si dice \textbf{lineare} (o \textbf{omomorfismo} di spazi vettoriali) se soddisfa le seguenti due proprietà:
    \begin{enumerate}
        \item \textbf{Additività:} \( \forall u, v \in V, \quad T(u+v) = T(u) + T(v) \)
        (Conserva la somma di vettori).
        \item \textbf{Omogeneità:} \( \forall v \in V, \forall \lambda \in K, \quad T(\lambda v) = \lambda T(v) \)
        (Conserva il prodotto per scalare).
    \end{enumerate}
    
    In base alle proprietà di \(T\) e agli spazi \(V, W\), un omomorfismo può essere classificato come:
    \begin{itemize}
        \item \textbf{Isomorfismo:} se \(T\) è un omomorfismo biettivo (iniettivo e suriettivo).
        \item \textbf{Monomorfismo:} se \(T\) è un omomorfismo iniettivo.
        \item \textbf{Epimorfismo:} se \(T\) è un omomorfismo suriettivo.
        \item \textbf{Endomorfismo:} se lo spazio di partenza e di arrivo coincidono, \(V = W\).
        \item \textbf{Automorfismo:} se \(V = W\) e \(T\) è biettiva (è un endomorfismo biettivo).
    \end{itemize}
}

% \Esempio{}{
%     Siano \(V = \mathbb{R}^2\) e \(W = \mathbb{R}^3\) (entrambi su campo \(K = \mathbb{R}\)).
%     Consideriamo l'applicazione \(T: \mathbb{R}^2 \to \mathbb{R}^3\) definita come:
%     \[
%     T( (x_1, x_2) ) = (x_1 + x_2, 2x_1, x_2 - x_1)
%     \]
%     Dimostriamo che \(T\) è un'applicazione lineare verificando le due proprietà.
    
%     Siano \(u = (x_1, x_2)\) e \(v = (y_1, y_2)\) due vettori generici in \(\mathbb{R}^2\), e \(\lambda \in \mathbb{R}\).
    
%     \begin{enumerate}
%         \item \textbf{Additività (\(T(u+v) = T(u) + T(v)\)):}
        
%         Calcoliamo prima \(u+v = (x_1+y_1, x_2+y_2)\).
%         \[
%         T(u+v) = T( (x_1+y_1, x_2+y_2) ) = ( (x_1+y_1) + (x_2+y_2), 2(x_1+y_1), (x_2+y_2) - (x_1+y_1) )
%         \]
%         \[
%         T(u+v) = ( x_1+y_1+x_2+y_2, 2x_1+2y_1, x_2+y_2-x_1-y_1 )
%         \]
        
%         Calcoliamo ora \(T(u) + T(v)\):
%         \[
%         T(u) + T(v) = (x_1+x_2, 2x_1, x_2-x_1) + (y_1+y_2, 2y_1, y_2-y_1)
%         \]
%         \[
%         T(u) + T(v) = ( (x_1+x_2)+(y_1+y_2), 2x_1+2y_1, (x_2-x_1)+(y_2-y_1) )
%         \]
%         \[
%         T(u) + T(v) = ( x_1+x_2+y_1+y_2, 2x_1+2y_1, x_2-x_1+y_2-y_1 )
%         \]
%         Confrontando i due risultati (a meno dell'ordine degli addendi), vediamo che \(T(u+v) = T(u) + T(v)\).

%         \item \textbf{Omogeneità (\(T(\lambda v) = \lambda T(v)\)):}
        
%         Calcoliamo prima \(\lambda v = (\lambda y_1, \lambda y_2)\).
%         \[
%         T(\lambda v) = T( (\lambda y_1, \lambda y_2) ) = ( (\lambda y_1) + (\lambda y_2), 2(\lambda y_1), (\lambda y_2) - (\lambda y_1) )
%         \]
%         \[
%         T(\lambda v) = ( \lambda(y_1+y_2), \lambda(2y_1), \lambda(y_2-y_1) )
%         \]
        
%         Calcoliamo ora \(\lambda T(v)\):
%         \[
%         \lambda T(v) = \lambda (y_1+y_2, 2y_1, y_2-y_1)
%         \]
%         \[
%         \lambda T(v) = ( \lambda(y_1+y_2), \lambda(2y_1), \lambda(y_2-y_1) )
%         \]
%         I due risultati coincidono.
%     \end{enumerate}
%     Avendo verificato entrambe le proprietà, \(T\) è un'applicazione lineare.
% }


\BoxBlue{Osservazioni:}{Proprietà delle Applicazioni Lineari}{
    Sia \(T: V \to W\) un'applicazione lineare.
    
    \begin{itemize}
        \item \textbf{T manda il vettore nullo nel vettore nullo:}
        L'immagine del vettore nullo di \(V\) è il vettore nullo di \(W\):
        \[ T(\bar{0}_V) = \bar{0}_W \]
        \textit{Dimostrazione:}
        Sfruttiamo la proprietà \( \bar{0}_V = 0 \cdot v \) (dove \(0\) è lo scalare) e l'omogeneità di \(T\):
        \[
        T(\bar{0}_V) = T(0 \cdot \bar{0}_V) \stackrel{\text{omogeneità}}{=} 0 \cdot T(\bar{0}_V) = \bar{0}_W
        \]
        (L'ultimo passaggio è vero perché il prodotto di uno scalare zero per qualsiasi vettore di \(W\) dà il vettore nullo \(\bar{0}_W\)).

        \item \textbf{Le applicazioni lineari conservano le combinazioni lineari:}
        Per ogni \(h \geq 1\), \(\forall u_1, \dots, u_h \in V\) e \(\forall \lambda_1, \dots, \lambda_h \in K\), si ha:
        \[
        T(\lambda_1 u_1 + \dots + \lambda_h u_h) = \lambda_1 T(u_1) + \dots + \lambda_h T(u_h)
        \]
        \textit{Dimostrazione (per induzione sul numero \(h\) di vettori):}
        
        \textbf{Base (h=1):}
        Dobbiamo dimostrare \( T(\lambda_1 u_1) = \lambda_1 T(u_1) \).
        Questo è vero per definizione, essendo la proprietà di \textbf{omogeneità} (la seconda proprietà delle applicazioni lineari).
        
        \textbf{Passo induttivo (h-1 \(\implies\) h):}
        Supponiamo (Ipotesi Induttiva) che la tesi sia vera per \(h-1\) vettori, cioè:
        \[
        T(\lambda_1 u_1 + \dots + \lambda_{h-1} u_{h-1}) = \lambda_1 T(u_1) + \dots + \lambda_{h-1} T(u_{h-1})
        \]
        Vogliamo dimostrare la tesi per \(h\) vettori. Consideriamo la combinazione lineare:
        \[
        T(\lambda_1 u_1 + \dots + \lambda_{h-1} u_{h-1} + \lambda_h u_h)
        \]
        Possiamo raggruppare i termini e usare le proprietà di \(T\). Sia \(u = \lambda_1 u_1 + \dots + \lambda_{h-1} u_{h-1}\) e \(v = \lambda_h u_h\).
        \[
        T(u + v) \stackrel{\text{additività}}{=} T(u) + T(v)
        \]
        Sostituendo \(u\) e \(v\):
        \[
        T(\lambda_1 u_1 + \dots + \lambda_{h-1} u_{h-1}) + T(\lambda_h u_h)
        \]
        Applichiamo ora l'ipotesi induttiva al primo termine e l'omogeneità (passo base) al secondo termine:
        \[
        \underbrace{(\lambda_1 T(u_1) + \dots + \lambda_{h-1} T(u_{h-1}))}_{\text{da Ip. Induttiva}} + \underbrace{\lambda_h T(u_h)}_{\text{da Omogeneità}}
        \]
        Che è esattamente la tesi per \(h\):
        \[
        \lambda_1 T(u_1) + \dots + \lambda_{h-1} T(u_{h-1}) + \lambda_h T(u_h)
        \]
        La proposizione è dimostrata.
    \end{itemize}
}
\newpage
\Proposizione{Le applicazioni lineari conservano la dipendenza lineare}{
    Sia \(T: V \to W\) un'applicazione lineare tra due K-spazi vettoriali \(V\) e \(W\).
    
    Se l'n-upla \((u_1, \dots, u_n)\) di vettori di \(V\) è \textbf{linearmente dipendente},
    allora la n-upla delle immagini \((T(u_1), \dots, T(u_n))\) è un'n-upla di vettori di \(W\) \textbf{linearmente dipendente}.

    \Dimostrazione{
        \textbf{Ipotesi:} L'n-upla \((u_1, \dots, u_n)\) è linearmente dipendente.
        Questo significa che esistono scalari \((\alpha_1, \dots, \alpha_n) \in K^n\), \textbf{non tutti nulli}, tali che:
        \[
        \alpha_1 u_1 + \dots + \alpha_n u_n = \bar{0}_V
        \]
        (dove \(\bar{0}_V\) è il vettore nullo di \(V\)).
        
        \textbf{Tesi:} L'n-upla \((T(u_1), \dots, T(u_n))\) è linearmente dipendente.
        Dobbiamo cioè dimostrare che esistono scalari \((\beta_1, \dots, \beta_n) \in K^n\), \textbf{non tutti nulli}, tali che:
        \[
        \beta_1 T(u_1) + \dots + \beta_n T(u_n) = \bar{0}_W
        \]
        (dove \(\bar{0}_W\) è il vettore nullo di \(W\)).
        
        \textbf{Svolgimento:}
        Applichiamo l'applicazione lineare \(T\) ad entrambi i membri dell'equazione che abbiamo per ipotesi:
        \[
        T( \alpha_1 u_1 + \dots + \alpha_n u_n ) = T(\bar{0}_V)
        \]
        
        Per la \textbf{linearità} di \(T\) (proprietà di conservare le combinazioni lineari), il primo membro diventa:
        \[
        T( \alpha_1 u_1 + \dots + \alpha_n u_n ) = \alpha_1 T(u_1) + \dots + \alpha_n T(u_n)
        \]
        
        Per una proprietà fondamentale delle applicazioni lineari, il secondo membro (l'immagine del vettore nullo) è:
        \[
        T(\bar{0}_V) = \bar{0}_W
        \]
        
        Uguagliando i due risultati, otteniamo:
        \[
        \alpha_1 T(u_1) + \dots + \alpha_n T(u_n) = \bar{0}_W
        \]
        
        Ora ci basta scegliere \( (\beta_1, \dots, \beta_n) = (\alpha_1, \dots, \alpha_n) \).
        Poiché per ipotesi gli scalari \(\alpha_i\) non erano tutti nulli, anche i \(\beta_i\) non sono tutti nulli.
        
        Abbiamo quindi trovato una combinazione lineare non banale dei vettori \(T(u_i)\) che dà il vettore nullo in \(W\). Questo dimostra la tesi.
    }
}
\bigskip \bigskip \bigskip

\BoxBlue{Osservazioni:}{Altre proprietà delle applicazioni lineari}{
    \begin{itemize}
        \item Se \(T: V \to W\) è un isomorfismo, allora la sua applicazione inversa \(T^{-1}: W \to V\) è anch'essa un'applicazione lineare e quindi è un isomorfismo.
        
        \item Siano \(T: V \to W\) e \(T': W \to U\) due applicazioni lineari.
        Allora l'applicazione composta \(T' \circ T: V \to U\) (definita da \((T' \circ T)(v) = T'(T(v))\)) è un'applicazione lineare.
    \end{itemize}
}

 
\Definizione{Kernel (o nucleo)}{
    Data un'applicazione lineare \( T : V \to W \), si definisce \textbf{nucleo} (o \textbf{kernel}) di \( T \) l'insieme:
    \[
        \ker (T) = \{\, u \in V \mid T(u) = \bar{0}_W \,\}.
    \]

    In altre parole, il kernel di \( T \) è l'insieme di tutti i vettori di \( V \) che vengono mandati nel vettore nullo di \( W \) tramite \( T \).
}

\Proposizione{Proprietà fondamentali delle applicazioni lineari}{
    Sia \( T : V \to W \) un'applicazione lineare. Allora valgono le seguenti proprietà:
    
    \begin{enumerate}
        \item \( T \) è \textbf{suriettiva} se e solo se l'immagine coincide con il codominio:
        \[
            \operatorname{Im}(T) = W.
        \]

        \item \( T \) è \textbf{iniettiva} se e solo se il nucleo è banale (contiene solo il vettore nullo):
        \[
            \ker(T) = \{\bar{0}_V\}.
        \]

        \item Se \( S \subseteq V \) è un sottoinsieme di generatori, allora l'immagine del sottospazio generato da \(S\) è uguale al sottospazio generato dalle immagini dei generatori:
        \[
            T(L(S)) = L(T(S)).
        \]
    \end{enumerate}

    \Dimostrazione{
    \textbf{Punto (ii) — iniettività e nucleo banale.}

    \textbf{(}\(\Rightarrow\)\textbf{)}  
    Supponiamo che \(T\) sia iniettiva. Poiché \(T\) è lineare, \(T(\underline{0}_V)=\underline{0}_W\).
    Sia ora \(u\in V\) con \(u\neq\underline{0}_V\). Per iniettività deve valere \(T(u)\neq T(\underline{0}_V)=\underline{0}_W\).
    Quindi non esiste vettore non nullo mandato in \(\underline{0}_W\), ovvero
    \[
        \ker(T)=\{\underline{0}_V\}.
    \]

    \medskip

    \textbf{(}\(\Leftarrow\)\textbf{)}  
    Supponiamo ora che \(\ker(T)=\{\underline{0}_V\}\) e mostriamo che \(T\) è iniettiva.
    Siano \(u,v\in V\) tali che \(T(u)=T(v)\). Allora
    \[
        T(u)-T(v)=\underline{0}_W \quad\Longrightarrow\quad T(u-v)=\underline{0}_W,
    \]
    quindi \(u-v\in\ker(T)=\{\underline{0}_V\}\). Di conseguenza \(u-v=\underline{0}_V\) e quindi \(u=v\).
    Questo dimostra l'iniettività di \(T\).

    \bigskip
    \textbf{Punto (iii) — \(T(L(S)) = L(T(S))\).}

    \textbf{Inclusione \(T(L(S)) \subseteq L(T(S))\).} \\
    Sia \(w\in T(L(S))\). Allora esiste \(u\in L(S)\) tale che \(T(u)=w\).
    Poiché \(u\in L(S)\), esistono vettori \(u_1,\dots,u_n\in S\) e scalari \(\alpha_1,\dots,\alpha_n\in K\) tali che
    \[
        u=\alpha_1 u_1 + \dots + \alpha_n u_n.
    \]
    Applicando \(T\) e usando la linearità otteniamo
    \[
        w = T(u) = \alpha_1 T(u_1) + \dots + \alpha_n T(u_n),
    \]
    cioè \(w\) è una combinazione lineare di vettori di \(T(S)\). Quindi \(w\in L(T(S))\).
    Abbiamo dimostrato \(T(L(S)) \subseteq L(T(S))\).

    \medskip

    \textbf{Inclusione \(L(T(S)) \subseteq T(L(S))\).} \\
    Sia \(z\in L(T(S))\). Allora esistono vettori \(w_1,\dots,w_m\in T(S)\) e scalari \(\beta_1,\dots,\beta_m\in K\) tali che
    \[
        z = \beta_1 w_1 + \dots + \beta_m w_m.
    \]
    Per definizione di \(T(S)\), per ogni \(w_i\) esiste \(v_i\in S\) con \(w_i=T(v_i)\). Quindi
    \[
        z = \beta_1 T(v_1) + \dots + \beta_m T(v_m) = T(\beta_1 v_1 + \dots + \beta_m v_m).
    \]
    Poiché \(\beta_1 v_1 + \dots + \beta_m v_m \in L(S)\), ne segue che \(z \in T(L(S))\).
    Abbiamo dunque \(L(T(S)) \subseteq T(L(S))\).

    \medskip

    Combinando le due inclusioni si ottiene l'uguaglianza:
    \[
        T(L(S)) = L(T(S)).
    \]
    \(\Box\)
}
}


 
\Proposizione{Il nucleo è un sottospazio vettoriale}{
    Sia \(T: V \to W\) un'applicazione lineare. Allora \(\ker(T)\) è un sottospazio vettoriale di \(V\).

    \Dimostrazione{
    Per prima cosa osserviamo che \(T\) è lineare, quindi \(T(\underline{0}_V)=\underline{0}_W\).  
    Dunque \(\underline{0}_V \in \ker(T)\), cioè \(\ker(T)\) non è vuoto.

    \medskip

    Siano ora \(u,v \in \ker(T)\). Per definizione \(T(u)=\underline{0}_W\) e \(T(v)=\underline{0}_W\).  
    Mostriamo che \(u+v \in \ker(T)\) e \(\alpha u \in \ker(T)\) per ogni \(\alpha\in K\).

    \textbf{Chiusura per addizione.}  
    \[
        T(u+v) \overset{(1)}{=} T(u) + T(v) = \underline{0}_W + \underline{0}_W = \underline{0}_W,
    \]
    quindi \(u+v \in \ker(T)\).

    \medskip

    \textbf{Chiusura per moltiplicazione per scalare.}  
    Sia \(\alpha\in K\). Poiché \(T(u)=\underline{0}_W\),
    \[
        T(\alpha u) \overset{(2)}{=} \alpha\,T(u) = \alpha\,\underline{0}_W = \underline{0}_W,
    \]
    quindi \(\alpha u \in \ker(T)\).

    \medskip

    Avendo verificato che \(\ker(T)\) è non vuoto, chiuso rispetto all'addizione e chiuso rispetto alla moltiplicazione per scalare, concludiamo che \(\ker(T)\) è un sottospazio vettoriale di \(V\).
    \(\Box\)
}
}

\bigskip \bigskip
 
\Proposizione{Conservazione della linearità dell’indipendenza}{
    Sia \(T : V \to W\) un'applicazione lineare iniettiva.  
    Se \(u_1, \dots, u_n\) è una \(n\)-upla di vettori linearmente indipendenti di \(V\),  
    allora \(T(u_1), \dots, T(u_n)\) è una \(n\)-upla di vettori linearmente indipendenti di \(W\).

    \Dimostrazione{
    Dobbiamo mostrare che se \((\alpha_1, \dots, \alpha_n) \in K^n\) è tale che
    \[
        \alpha_1 T(u_1) + \dots + \alpha_n T(u_n) = \underline{0}_W,
    \]
    allora necessariamente \(\alpha_1 = \dots = \alpha_n = 0\).

    \medskip

    Per linearità di \(T\) si ha:
    \[
        T(\alpha_1 u_1 + \dots + \alpha_n u_n) = \underline{0}_W.
    \]
    Questo implica che
    \[
        \alpha_1 u_1 + \dots + \alpha_n u_n \in \ker(T).
    \]
    Poiché \(T\) è iniettiva, \(\ker(T) = \{\underline{0}_V\}\).  
    Dunque:
    \[
        \alpha_1 u_1 + \dots + \alpha_n u_n = \underline{0}_V.
    \]

    \medskip

    Ma i vettori \(u_1, \dots, u_n\) sono linearmente indipendenti, quindi
    \[
        \alpha_1 = \alpha_2 = \dots = \alpha_n = 0.
    \]

    \medskip

    Pertanto, \(T(u_1), \dots, T(u_n)\) sono linearmente indipendenti in \(W\).
    \(\Box\)
}
}


\newpage
\Teorema{Teorema dell'Equazione Dimensionale}{
    Sia \(T: V \to W\) un'applicazione lineare. Supponiamo che \(V\) sia uno spazio vettoriale finitamente generato, con \(\dim V = n\).
    
    Allora \(\ker(T)\) e \(\operatorname{Im}(T)\) sono anch'essi finitamente generati e vale la seguente relazione:
    \[
    \dim V = \dim(\ker(T)) + \dim(\operatorname{Im}(T))
    \]
    (La dimensione dello spazio di partenza è la somma della dimensione del nucleo e della dimensione dell'immagine).
}

\BoxBlue{Osservazioni:}{conseguenze del Teorema dell'equazione Dimensionale}{
    Sia \(T: V \to W\) un'applicazione lineare, con \(\dim V = n\) e \(\dim W = m\).
    \begin{itemize}
        \item Se \(n > m\) (\(\dim V > \dim W\)), \(T\) \textbf{non può essere iniettiva}.
        \newline
        \textit{Dim:} Sappiamo \(\dim(\operatorname{Im}(T)) \leq \dim W = m\).
        Dalla formula: \(\dim(\ker(T)) = n - \dim(\operatorname{Im}(T)) \geq n - m > 0\).
        Poiché \(\dim(\ker(T)) > 0\), il nucleo non è banale e \(T\) non è iniettiva.
        
        \item Se \(n < m\) (\(\dim V < \dim W\)), \(T\) \textbf{non può essere suriettiva}.
        \newline
        \textit{Dim:} Sappiamo \(\dim(\ker(T)) \geq 0\).
        Dalla formula: \(\dim(\operatorname{Im}(T)) = n - \dim(\ker(T)) \leq n\).
        Poiché \(\dim(\operatorname{Im}(T)) \leq n < m\), l'immagine non può coincidere con \(W\).
        
        \item Se \(n = m\) (\(\dim V = \dim W\)), \(T\) è \textbf{iniettiva} \(\iff\) \(T\) è \textbf{suriettiva}.
        \newline
        \textit{Dimostrazione:} Sia \(\dim V = \dim W = n\).
        Dall'equazione dimensionale sappiamo che \(n = \dim(\ker(T)) + \dim(\operatorname{Im}(T))\).
        \begin{itemize}
            \item[\((\Rightarrow)\)] \textbf{(Iniettiva \(\implies\) Suriettiva)}
            Se \(T\) è iniettiva, allora \(\dim(\ker(T)) = 0\).
            Sostituendo nell'equazione:
            \[ n = 0 + \dim(\operatorname{Im}(T)) \implies \dim(\operatorname{Im}(T)) = n \]
            L'immagine \(\operatorname{Im}(T)\) è un sottospazio di \(W\) che ha la stessa dimensione di \(W\) (\(n\)). Per la proprietà dei sottospazi (punto 3 della prop. precedente), questo implica che \(\operatorname{Im}(T) = W\).
            Quindi \(T\) è suriettiva.
            
            \item[\((\Leftarrow)\)] \textbf{(Suriettiva \(\implies\) Iniettiva)}
            Se \(T\) è suriettiva, allora \(\operatorname{Im}(T) = W\), e quindi \(\dim(\operatorname{Im}(T)) = \dim W = n\).
            Sostituendo nell'equazione:
            \[ n = \dim(\ker(T)) + n \implies \dim(\ker(T)) = 0 \]
            Poiché \(\dim(\ker(T)) = 0\), il nucleo è banale (\(\ker(T) = \{\bar{0}_V\}\)).
            Quindi \(T\) è iniettiva.
        \end{itemize}
        Di conseguenza, quando le dimensioni coincidono, T è iniettiva se e solo se è suriettiva, e quindi se e solo se è un \textbf{isomorfismo}.
    \end{itemize}
}
 


 
